\documentclass[12pt, a4paper, oneside]{ctexart}
\usepackage{amsmath, amsthm, amssymb, bm, xcolor, tikz, framed, graphicx, hyperref, mathrsfs}
\usepackage{bookmark}
\usepackage{pgfplots}
\pgfplotsset{compat=1.18}


\title{Quantum Mechanics Homework (Week 3)}
\author{Tang Yi \\ [学号已脱敏]}
\usepackage[a4paper, left=2cm, right=2cm, top=2cm, bottom=2cm]{geometry}
\linespread{1.3} % 行距更紧凑
\definecolor{shadecolor}{RGB}{241, 241, 255}

\newcounter{problemname}
\newenvironment{problem}
  {\begin{shaded}\stepcounter{problemname}\par\noindent\textbf{题目 \theproblemname. }}
  {\end{shaded}\par}
\newenvironment{solution}
  {\par\noindent\textbf{题目 \theproblemname 的解答. }}
  {\par}
\newenvironment{note}
  {\par\noindent\textbf{题目 \theproblemname 的注记. }}
  {\par}

\begin{document}

\maketitle

\begin{problem}
    对一维无限深方势阱的第$n$个定态分别计算$\langle x \rangle ,\langle p \rangle , \Delta_x $和$ \Delta_p$ 。说明何态最接近量子极限，何态最接近经典极限。
\end{problem}

\begin{solution}
    不难发现，一维无限深方势阱的第$n$个定态的波函数为

    \centerline{$\psi_n(x) = \sqrt{\dfrac{2}{a}} \sin\left(\dfrac{n\pi x}{a}\right)$}

    则

    \centerline{$\langle x \rangle_n = \displaystyle\int_0^a \mathrm{d}x\,x\psi^*\psi = \dfrac{a}{2}$}

    \centerline{$\langle p \rangle_n = \displaystyle\int_0^a \mathrm{d}x\,\psi^*\left(-i\hbar\dfrac{\partial}{\partial x}\right)\psi = 0$}

    \centerline{$\langle x^2 \rangle_n = \displaystyle\int_0^a \mathrm{d}x\,x^2\psi^*\psi = \dfrac{a^2}{3} - \dfrac{a^2}{2n^2\pi^2}$}

    \centerline{$\langle p^2 \rangle_n = \displaystyle\int_0^a \mathrm{d}x\,\psi^*\left(-\hbar^2\dfrac{\partial^2}{\partial x^2}\right)\psi = \dfrac{n^2\pi^2\hbar^2}{a^2}$}

    因此

    \centerline{$\Delta x_n = \sqrt{\langle x^2 \rangle_n - \langle x \rangle_n^2} = \dfrac{a}{2\pi n}\sqrt{\dfrac{1}{3}n^2\pi^2 - 2}$}

    \centerline{$\Delta p_n = \sqrt{\langle p^2 \rangle_n - \langle p \rangle_n^2} = \dfrac{n\pi\hbar}{a}$}

    由此可见，$\Delta x_n \Delta p_n = \dfrac{\hbar}{2}\sqrt{\dfrac{1}{3}n^2\pi^2 - 2} $，当$n$较小的时候，其与量子极限$\dfrac{\hbar}{2}$相差不大，因此较接近量子极限。

    当$n$较大的时候，$\dfrac{\Delta E_n}{E_n} \sim \dfrac{1}{n}$，能量的间隙较小，因此较接近经典极限。
    
\end{solution}

\begin{problem}
  设处于一维无限深方势阱中的粒子初始时刻的波函数为: $\varPsi(x,0) =Ax(x−a),0 < x < a $。

  (1) 求归一化系数$A$；

  (2) 求测得粒子处于能量本征态$\psi_n(x)$的几率$|c_n|^2$ ；

  (3) 写出$t$时刻的波函数；

  (4) 求$t$时刻能量的期望值。
\end{problem}

\begin{solution}
    （1）归一化条件要求

    \begin{align*}
      1 &= \langle \varPsi | \varPsi \rangle \\
      &= A^2 \displaystyle\int_0^a \mathrm{d}x\,x^2(x-a)^2 \\
      &= A^2 \left[\dfrac{a^5}{5} - \dfrac{2a^5}{4} + \dfrac{a^5}{3}\right] \\
      &= \dfrac{1}{30} A^2 a^5
    \end{align*}

    因此，$A = \sqrt{\dfrac{30}{a^5}}$

    （2）能量本征态为$\psi_n(x) = \sqrt{\dfrac{2}{a}}\sin\left(\dfrac{n\pi x}{a}\right)$，因此

        对$c_n$，有

        \begin{align*}
          c_n &= \langle \psi_n | \varPsi \rangle \\
          &= \displaystyle\int_0^a \mathrm{d}x\,\psi_n^*(x)\varPsi(x,0) \\
          &= \sqrt{\dfrac{60}{a^6}}\left( \displaystyle\int_0^a \mathrm{d}x x^2 \sin\left(\dfrac{n \pi x}{a}\right) - a \displaystyle\int_0^a \mathrm{d}x x \sin\left(\dfrac{n \pi x}{a}\right)\right) \\
          &= \sqrt{\dfrac{60}{a^6}} \dfrac{2a^3}{n^3 \pi^3} \left(\cos(n \pi) - 1 \right)\\
          &= - \dfrac{4 \sqrt{15}}{n^3 \pi^3} \left(1 - (-1)^n \right)\\
          &= \begin{cases}
            - \dfrac{8 \sqrt{15}}{n^3 \pi^3}, & n \,\text{is an odd number;} \\
            0, & n \,\text{is a even number.}
          \end{cases}
        \end{align*}

        则$|c_n|^2$是

        \centerline{$|c_n|^2 = \begin{cases}
          \dfrac{960}{n^6 \pi^6}, & n \,\text{is an odd number;} \\
          0, & n \,\text{is a even number.}
        \end{cases}$}

    （3）$t$时刻的波函数为

    \begin{align*}
      \varPsi(x,t) &= \displaystyle\sum_{n=1}^{\infty} c_n \psi_n(x) e^{-iE_nt/\hbar} \\
      &= \displaystyle\sum_{\substack{n=1 \\ n\,\text{is an odd number}}}^{\infty} - \dfrac{8 \sqrt{15}}{n^3 \pi^3} \sqrt{\dfrac{2}{a}}\sin\left(\dfrac{n\pi x}{a}\right) \exp\left[{-i\dfrac{n^2 \pi^2 \hbar}{2ma^2}t}\right] \\
      &= - \dfrac{8}{\pi^3} \sqrt{\dfrac{30}{a}} \displaystyle \sum_{k=0}^{\infty} \dfrac{1}{(2k+1)^3} \sin\left(\dfrac{(2k+1)\pi x}{a}\right) \exp\left[{-i\dfrac{(2k+1)^2 \pi^2 \hbar}{2ma^2}t}\right]
    \end{align*}

    （4）能量的期望值为

    \begin{align*}
      \langle E \rangle &= \langle \varPsi | \hat{H} | \varPsi \rangle \\
      &= \left\langle \left.\displaystyle\sum_m c_m \psi_m e^{-iE_mt/\hbar} \right| \hat{H} \left| \displaystyle\sum_n c_n \psi_n e^{-iE_nt/\hbar} \right.\right\rangle \\
      &= \displaystyle\sum_{m} \displaystyle\sum_n c_m^* c_n e^{i(E_m - E_n)t/\hbar} \langle \psi_m | \hat{H} | \psi_n \rangle \\
      &= \displaystyle\sum_{m} \displaystyle\sum_n c_m^* c_n e^{i(E_m - E_n)t/\hbar} E_n \delta_{mn} \\
      &= \displaystyle\sum_n |c_n|^2 E_n \\
      &= \displaystyle\sum_{\substack{n=1 \\ n\,\text{is an odd number}}}^{\infty} \dfrac{960}{n^6 \pi^6} \dfrac{n^2 \pi^2 \hbar^2}{2ma^2} \\
      &= \dfrac{480 \hbar^2}{ma^2 \pi^4} \displaystyle\sum_{k=0}^{\infty} \dfrac{1}{(2k+1)^4} \\
    \end{align*}

    利用mma，我们不难计算得

    \centerline{$\displaystyle\sum_{k=0}^{\infty} \dfrac{1}{(2k+1)^4} = \dfrac{\pi^4}{96}$}

    代入，则

    \centerline{$\langle E \rangle = \dfrac{5 \hbar^2}{ma^2}$}
\end{solution}

\begin{problem}
    设处于一维无限深方势阱中的粒子初始时刻的波函数由前两个能量本征态叠加而成：

    \centerline{$\varPsi(x,0) = A[\psi_1(x)+2\psi_2(x)]$}

    (1) 求归一化系数$A$；

    (2) 求$\varPsi(x,t)$和 $|\varPsi(x,t)|^2$；

    (3) 求$t$时刻位置和动量的期望值。

\end{problem}

\begin{solution}
  （1）归一化条件要求
  \begin{align*}
    1 &= \langle \varPsi | \varPsi \rangle \\
    &= A^*A \left(\langle \psi_1 | \psi_1 \rangle + 4\langle \psi_2 | \psi_2 \rangle + 2\langle \psi_1 | \psi_2 \rangle + 2\langle \psi_2 | \psi_1 \rangle\right) \\
    &= 5 A^*A
  \end{align*}

  为了方便，我们把$A$取为实数，因此，$A = \dfrac{1}{\sqrt{5}}$。

  （2）$t$时刻的波函数为：

  \centerline{$\varPsi(x,t) = \dfrac{1}{\sqrt{5}}\psi_1(x)e^{-iE_1t/\hbar} + \dfrac{2}{\sqrt{5}}\psi_2(x)e^{-iE_2t/\hbar}$}


  我们设$(\psi_1,\psi_2) = \psi_1^* \psi_2$，则

  概率密度函数：

  \begin{align*}
    |\Psi(x,t)|^2 &= \left( \varPsi , \varPsi \right)\\
    &= \left( \dfrac{1}{\sqrt{5}}\psi_1(x)e^{-iE_1t/\hbar} + \dfrac{2}{\sqrt{5}}\psi_2(x)e^{-iE_2t/\hbar}, \dfrac{1}{\sqrt{5}}\psi_1(x)e^{-iE_1t/\hbar} + \dfrac{2}{\sqrt{5}}\psi_2(x)e^{-iE_2t/\hbar}\right) \\
    &= \dfrac{1}{5} \left[ ( \psi_1 , \psi_1 ) + 2 ( \psi_1,\psi_2 ) e^{i(E_1 - E_2)t / \hbar} + 2 ( \psi_2,\psi_1 ) e^{i(E_2 - E_1)t / \hbar} + 4 ( \psi_2 ,\psi_2 )\right]
  \end{align*}

  由于$\psi_1,\psi_2$都是实的，因此

  \centerline{$( \psi_1 , \psi_2 ) = ( \psi_1 , \psi_2 ) ^* = ( \psi_2 , \psi_1 ) $}

  代入上式，得

  \begin{align*}
    |\Psi(x,t)|^2 &= \dfrac{1}{5} \left[ ( \psi_1 , \psi_1 ) + 4 ( \psi_2 ,\psi_2 ) + 4 ( \psi_1,\psi_2 ) \cos\left(\dfrac{E_2 - E_1}{\hbar}t\right) \right] \\
    &= \dfrac{2}{5a} \left[  \sin^2\left(\dfrac{\pi x}{a}\right) + 4\sin^2\left(\dfrac{2\pi x}{a}\right) + 4 \sin\left(\dfrac{\pi x}{a}\right)\sin\left(\dfrac{2\pi x}{a}\right) \cos\left(\dfrac{3\pi^2\hbar}{2ma^2}t\right) \right]
  \end{align*}

  （3）先算位置
  \begin{align*}
    \langle x \rangle &= \langle \varPsi | x | \varPsi \rangle \\
    &= \left\langle \left.\dfrac{1}{\sqrt{5}}\psi_1 e^{-iE_1t/\hbar} + \dfrac{2}{\sqrt{5}}\psi_2 e^{-iE_2t/\hbar} \right| x \left| \dfrac{1}{\sqrt{5}}\psi_1 e^{-iE_1t/\hbar} + \dfrac{2}{\sqrt{5}}\psi_2 e^{-iE_2t/\hbar} \right.\right\rangle \\
    &= \dfrac{1}{5} \left[ \langle \psi_1 | x|\psi_1 \rangle + 4\langle \psi_2 | x|\psi_2 \rangle +2\langle \psi_1 | x|\psi_2 \rangle e^{i(E_1 - E_2)t/\hbar} + 2\langle\psi_2|x|\psi_1\rangle e^{i(E_2 - E_1)t/\hbar}\right] \\
  \end{align*}
  
  同时，由于$x$是厄米的，且波函数是实的，因此
  \begin{align*}
    \langle \psi_1 | x|\psi_2 \rangle &= \langle \psi_1 | x|\psi_2 \rangle^* \\
    &= \langle \psi_2 | x^\dagger|\psi_1 \rangle\\
    &= \langle \psi_2 | x|\psi_1 \rangle
  \end{align*}
  
  同时，可以计算
  \begin{align*}
    \langle \psi_1 | x|\psi_1 \rangle &= \dfrac{a}{2} \\
    \langle \psi_2 | x|\psi_2 \rangle &= \dfrac{a}{2} \\
    \langle \psi_1 | x|\psi_2 \rangle &= -\dfrac{16a}{9\pi^2}
  \end{align*}
  
  代入上式，得

  \centerline{$\langle x \rangle = \dfrac{a}{2} - \dfrac{64a}{45\pi^2}\cos\left(\dfrac{3\pi^2\hbar}{2ma^2}t\right)$}

  再代入式子

  \centerline{$\dfrac{\mathrm{d} \langle x \rangle}{\mathrm{d}t} = \dfrac{\langle p \rangle}{m}$}

  可得

  \centerline{$\langle p \rangle = \dfrac{32\hbar}{15a}\sin\left(\dfrac{3\pi^2\hbar}{2ma^2}t\right)$}
\end{solution}

\begin{problem}
  考虑一维情况，粒子在一般的势场中运动，$\psi_n(x),\psi_m(x)$ 是分别属于$E_m,E_n$ 的束缚态本征函数，且$E_n \neq E_m$。从定态薛定谔方程出发，证明：

  \centerline{$\displaystyle\int_{-\infty}^{\infty} \psi_n(x) \psi_m(x)\mathrm{d}x = 0$}
\end{problem}

\begin{solution}
  由Schrödinger方程：

  \centerline{$\hat{H}|\psi_m\rangle = E_m |\psi_m\rangle,\hat{H}|\psi_n\rangle = E_n |\psi_n\rangle$}

  不难有：

  \begin{align*}
    \langle \psi_n | \hat{H} | \psi_m \rangle &= E_m \langle \psi_n | \psi_m \rangle \\
    &= E_m \langle \psi_m | \psi_n \rangle^* \\
    &= \langle \psi_m | \hat{H} | \psi_n \rangle^* \\
    &= E_n \langle \psi_m | \psi_n \rangle^*\\
    &= E_n \langle \psi_n | \psi_m \rangle
  \end{align*}

  因此有

  \centerline{$(E_m - E_n) \langle \psi_n | \psi_m \rangle = 0$}

  又因为$E_m \neq E_n$，因此

  \centerline{$\langle \psi_n | \psi_m \rangle = 0$}

  得证！
\end{solution}

\begin{problem}
  设粒子处于半壁无限高的势场中
  \[
    V(x) = \begin{cases}
      \infty, & x < 0 \\
      -V_0, & 0<x <a\\
      0, & x > a
    \end{cases}
  \]
  求至少存在一个束缚态的条件。（提示：束缚态的能量$E < 0$）
\end{problem}

\begin{solution}
  在$0<x<a$范围内,有薛定谔方程:

  \centerline{$-\dfrac{\hbar^2}{2m}\partial_x^2 \psi_1 -V_0 \psi_1 = E \psi_1$}

  化简后,有

  \centerline{$\partial_x^2 \psi_1 + k_1^2 \psi_1 = 0$}

  其中,$k_1 = \sqrt{\dfrac{2m(E+V_0)}{\hbar^2}}$.

  在$x>a$范围内,有薛定谔方程:

  \centerline{$-\dfrac{\hbar^2}{2m}\partial_x^2 \psi_2 = E \psi_2$}

  化简后,有

  \centerline{$\partial_x^2 \psi_2 - k_2^2 \psi_2 = 0$}

  其中,$k_2 = \sqrt{\dfrac{-2mE}{\hbar^2}}$.

  同时,由于连续性条件,以及无穷远收敛条件,应该有
  \begin{align*}
    \psi_1(0) &= 0 \\
    \psi_1(a) &= \psi_2(a) \\
    \partial_x \psi_1(a) &= \partial_x \psi_2(a)\\
    \lim_{x \to +\infty} \psi_2(x) &= 0
  \end{align*}

  因此,不难有

  \centerline{$\psi_1(x) = A \sin(k_1 x)$}

  \centerline{$\psi_2(x) = B e^{-k_2 x}$}

  代入边界条件,有

  \begin{align*}
    A \sin(k_1 a) &= B e^{-k_2 a} \\
    A k_1 \cos(k_1 a) &= -B k_2 e^{-k_2 a}
  \end{align*}

  消去$A,B$,代入$k_1,k_2$,有

  \centerline{$\cot\left(a\sqrt{\dfrac{2m(E+V_0)}{\hbar^2}}\right) = -\sqrt{-\dfrac{E}{E+V_0}}$}

  设$\xi = a\sqrt{\dfrac{2m(E+V_0)}{\hbar^2}}$,则

  \centerline{$\cot \xi = - \sqrt{\dfrac{\dfrac{2mV_0a^2}{\hbar^2}}{\xi^2} - 1}$}

  上式说明,$\cot \xi < 0 $,即$\dfrac{\pi}{2}+k\pi < \xi < \pi+k\pi,k\in\mathbb{Z}$.

  两边平方化简后,利用三角恒等式有

  \centerline{$\xi^2 \csc^2 \xi = \dfrac{2mV_0a^2}{\hbar^2}$}

  由于$E<0$,因此$0 < \xi < a\sqrt{\dfrac{2mV_0}{\hbar^2}}$.

  因此,两边开方化简有

  \centerline{$\sin \xi = \dfrac{\hbar}{\sqrt{2mV_0}a} \xi$}

  要使其有解,显然需要

  \centerline{$\dfrac{\hbar}{\sqrt{2mV_0}a} < 1$}

  即

  \centerline{$V_0 > \dfrac{\hbar^2}{2ma^2}$}

  除此之外,还需要让上面计算的两个$\xi$范围有交集,则至少需要

  \centerline{$a\sqrt{\dfrac{2mV_0}{\hbar^2}} > \dfrac{\pi}{2}$}

  即

  \centerline{$V_0 > \dfrac{\pi^2 \hbar^2}{8ma^2}$}

  取交集,综上,需要有束缚态,需要满足

  \centerline{$V_0 > \dfrac{\pi^2 \hbar^2}{8ma^2}$}
\end{solution}

\begin{problem}
  考虑双$\delta$函数势阱：$V(x) = -\gamma [\delta(x+a)+\delta(x-a)]$，其中$\gamma$和$a$都是正的常数。

  （1）写出束缚态本征能量满足的方程；

  （2）基于第一问的结果分别讨论$a \to 0$和$a \to \infty$时束缚态解的情况。
\end{problem}

\begin{solution}
  (1) 显然满足薛定谔方程

  \centerline{$-\dfrac{\hbar^2}{2m}\partial_x^2 \psi - \gamma [\delta(x+a)+\delta(x-a)]\psi = E \psi$}

  在$x \neq \pm a$时,有薛定谔方程

  \centerline{$\partial_x^2 \psi + \dfrac{2mE}{\hbar^2} \psi = 0$}

  因此,解的形式为

  \centerline{$\psi = A e^{kx}+B e^{-kx}$}

  其中,$k = \sqrt{\dfrac{-2mE}{\hbar^2}}$.

  同时,波函数需要满足连续性条件,无穷远收敛,宇称对称性.因此,结果只可能是

  奇宇称:
  \[
  \psi(x) = 
  \begin{cases}
  C e^{k(x+a)}, & x < -a \\[1ex]
  - C \dfrac{e^{kx}-e^{-kx}}{e^{ka}-e^{-ka}}, & -a < x < a \\[2ex]
  - C e^{-k(x-a)}, & x > a
  \end{cases}
  \]

  偶宇称:

  \[
  \psi(x) = 
  \begin{cases}
  D e^{k(x+a)}, & x < -a \\[1ex]
  D \dfrac{e^{kx}+e^{-kx}}{e^{ka}+e^{-ka}}, & -a < x < a \\[2ex]
  D e^{-k(x-a)}, & x > a
  \end{cases}
  \]

  下面我们带入$\delta$函数的条件,即

  \centerline{$\partial_x \psi|_{x=-a^+} - \partial_x \psi|_{x=-a^-} = -\dfrac{2m\gamma}{\hbar^2}\psi(-a)$}
  \centerline{$\partial_x \psi|_{x=a^+} - \partial_x \psi|_{x=a^-} = -\dfrac{2m\gamma}{\hbar^2}\psi(a)$}

  对于奇宇称,有

  \centerline{$k = \dfrac{m\gamma}{\hbar^2}(1 - e^{-2ka})$}

  对于偶宇称,有

  \centerline{$k = \dfrac{m\gamma}{\hbar^2}(1 + e^{-2ka})$}

  其中$k$是$E$的单值函数.
  
  \newpage

  (2) 当$a \to 0$时:
  
  奇宇称的束缚态能量$E=-\dfrac{\hbar^2k^2}{2m}=-\dfrac{m\gamma^2}{2\hbar^2}(1-e^{-2ka})^2 \to 0$,这显然对应束缚态$E<0$,因此此时奇宇称消失.

  偶宇称的束缚态能量$E=-\dfrac{\hbar^2k^2}{2m}=-\dfrac{m\gamma^2}{2\hbar^2}(1+e^{-2ka})^2 \to -\dfrac{2m\gamma^2}{\hbar^2}$,这显然对应束缚态$E<0$,因此此时偶宇称存在,退化为中心单势阱.

  当$a \to \infty$时:

  奇宇称的束缚态能量$E=-\dfrac{\hbar^2k^2}{2m}=-\dfrac{m\gamma^2}{2\hbar^2}(1-e^{-2ka})^2 \to -\dfrac{m\gamma^2}{2\hbar^2}$,这显然对应束缚态$E<0$,因此此时奇宇称存在.

  偶宇称的束缚态能量$E=-\dfrac{\hbar^2k^2}{2m}=-\dfrac{m\gamma^2}{2\hbar^2}(1+e^{-2ka})^2 \to -\dfrac{m\gamma^2}{2\hbar^2}$,这显然对应束缚态$E<0$,因此此时偶宇称存在.

  此时能量都与单势阱相同,这是因为两势阱足够远导致解耦合.
\end{solution}
\end{document}